\label{p3-20260826:chap:medida-accion-positiva}
\index{medida!catálogo de emisiones}
\index{acción!cociclo positivo}
\index{traza normalizada}

La sección anterior identifica isométricamente un sector positivo de rango tres en
la interfaz orientada \(t=7\), la dodecafase y el módulo de incidencia de Witt. La construcción
siguiente deriva el otro coeficiente de la ley angular, \(5/468\), y
demuestra por qué ambos términos se combinan mediante una suma positiva. La
cuestión exige separar
tres operaciones:
\begin{enumerate}
  \item contar presentaciones microscópicas;
  \item contar emisiones distintas después de formar el cociente TPK;
  \item asignar una acción aditiva a una ruta.
\end{enumerate}
Las dos primeras operaciones producen medidas diferentes. La tercera se
caracteriza por una propiedad universal y no por la proximidad posterior de
un decimal.

Denotaremos por \(\Gtr\) el grafo dirigido de \(1944\) estados del lector de
caminos enriquecidos. Su categoría libre de caminos se escribirá
\(\mathsf P(\Gtr)\). La definición local de sus aristas y de su acción se
desarrolla en la sección correspondiente.

\section{Espacio de semillas, cociente de emisiones y fibras regionales}

Sea
\[
 \Omega_{\rm seed}
 =\bigl((\mathbb Z/9\mathbb Z)^2\times D_4\bigr)^2,
 \qquad
 |\Omega_{\rm seed}|=104\,976,
\]
el conjunto ejecutable de parejas de semillas del emisor finito, y sea
\[
 F:\Omega_{\rm seed}\twoheadrightarrow\mathcal U_6,
 \qquad
 |\mathcal U_6|=468,
\]
la aplicación que conserva la emisión séxtuple distinta. El censo exacto de
fibras es
\[
\begin{array}{c|c|c}
|F^{-1}(U)|&\#\{U\text{ con ese tamaño}\}&\text{semillas aportadas}\\
\hline
144&432&62\,208\\
432&18&7\,776\\
1\,944&18&34\,992.
\end{array}
\]
Las tres contribuciones suman \(104\,976\). Por tanto, \(F\) no es una
cubierta regular: emisiones diferentes poseen números diferentes de
presentaciones microscópicas.

La microfibra de \(\pi\) contiene cinco emisiones distintas. Cuatro tienen
multiplicidad \(144\) y una tiene multiplicidad \(432\). Este dato basta
para demostrar que la medida uniforme sobre semillas y la medida uniforme
sobre emisiones no son la misma.

\section{\(2\) medidas exactas sobre el cociente de emisiones y palabras}

Si cada semilla recibe peso \(1/104\,976\), la medida imagen por \(F\) es
\[
 \nu_{\rm seed}(U)
 =\frac{|F^{-1}(U)|}{104\,976}.
\]
Sus átomos posibles son
\[
 \frac1{729},\qquad\frac1{243},\qquad\frac1{54}.
\]
\Needspace{5\baselineskip}
En particular,
\[
 \nu_{\rm seed}(\mathscr R_\pi^{\rm micro})
 =\frac{4\cdot144+432}{104\,976}
 =\boxed{\frac7{729}}.
\]

Si el estado elemental es, en cambio, una emisión distinta
\(U\in\mathcal U_6\), el álgebra diagonal de observables es
\[
 \mathcal D_{468}=\ell^\infty(\mathcal U_6)
\]
y su estado de conteo normalizado es
\[
 \tau_{468}(f)
 =\frac1{468}\sum_{U\in\mathcal U_6}f(U).
\]
Sea \(\Pi_\pi^{(468)}\) la función característica de las cinco regiones de la
microfibra. Entonces
\[
 \boxed{\tau_{468}(\Pi_\pi^{(468)})=\frac5{468}.}
\]
La diferencia exacta entre ambas lecturas es
\[
 \frac5{468}-\frac7{729}
 =\frac{41}{37\,908}.
\]
No son dos aproximaciones. La primera medida cuenta presentaciones; la
segunda cuenta clases de emisión.

\begin{theorem}[Unicidad de la medida uniforme del catálogo]
\label{p3-20260826:thm:unicidad-medida-468}
La única probabilidad sobre \(\mathcal U_6\) invariante bajo toda
relabelización de sus \(468\) átomos es la medida uniforme.
Equivalentemente, la única traza positiva y unital sobre
\(M_{468}(\mathbb C)\) es
\[
 \tau_{468}(A)=\frac1{468}\operatorname{Tr}A.
\]
\end{theorem}

\begin{proof}
Las transposiciones llevan cualquier átomo a cualquier otro. La invariancia
fuerza que todos tengan la misma masa, y la normalización la fija en
\(1/468\). En la versión matricial, la invariancia unitaria iguala el valor
de todos los proyectores de rango uno. La descomposición espectral y la
linealidad reconstruyen la traza normalizada.
\end{proof}

La hipótesis decisiva está a la vista: el lector trabaja con el
\emph{catálogo cociente}. El teorema no afirma que la medida uniforme de
semillas se convierta en uniforme al olvidar las fibras. La simetría usada
es la naturalidad del objeto finito sin etiquetas, no una propiedad
ergódica que se atribuya sin prueba a la dinámica microscópica.

\section{Hilbertización canónica del cociente}

La relación entre semillas y emisiones puede formularse sin ambigüedad.
Sean
\[
 H_{\rm seed}=\ell^2(\Omega_{\rm seed}),
 \qquad
 H_U=\ell^2(\mathcal U_6).
\]
El emisor determina
\[
 \boxed{
 J_F\delta_U
 =|F^{-1}(U)|^{-1/2}
 \sum_{\omega\in F^{-1}(U)}\delta_\omega.
 }
\]
Las fibras distintas son disjuntas y cada columna posee norma uno; por
tanto,
\[
 J_F^*J_F=I_{H_U}.
\]
Además, para todo observable \(f\) del catálogo,
\[
 M_{f\circ F}J_F=J_FM_f.
\]
La imagen \(J_FH_U\) es exactamente el subespacio de funciones constantes
en cada fibra.

Si \(F^\sharp:H_U\to H_{\rm seed}\) denota la aplicación inducida por
precomposición, sin normalizar,
entonces
\[
 J_F
 =F^\sharp\bigl((F^\sharp)^*F^\sharp\bigr)^{-1/2}.
\]
Es el factor isométrico de la descomposición polar de esa aplicación. Esta
expresión demuestra su naturalidad bajo isomorfismos y relabelizaciones del
emisor. No implica una functorialidad falsa bajo cualquier composición de
sobreyeciones no regulares: la normalización depende de las fibras de la
aplicación concreta.

En lenguaje de medidas, la elevación de la traza de cociente que es
uniforme dentro de cada fibra vale
\[
 \widetilde\tau(\omega)
 =\frac1{468\,|F^{-1}(F\omega)|}.
\]
Cada fibra recibe masa total \(1/468\). La compensación inversa por
multiplicidad, y no una invocación informal a órbitas, es la prueba exacta.

\begin{corollary}[Canonicidad de \(5/468\)]
\label{p3-20260826:cor:canonicidad-cinco-468}
Una vez que el lector regional se define sobre emisiones distintas, la masa
de la microfibra de \(\pi\) queda forzada:
\[
 \tau_{468}(\Pi_\pi^{(468)})
 =\frac{\operatorname{rank}\Pi_\pi^{(468)}}{468}
 =\frac5{468}.
\]
\end{corollary}

\section{Defecto positivo entre las medidas regional y dodecafásica}

Escribamos
\[
 H_{\rm ang}=H_U\otimes V_{11}.
\]
El escalar interno
\[
 \Delta_4
 =\frac{\pi-e\log\pi}{270}
 \left(1+\frac\pi{729}\right)
 =0.0001111964226921402\ldots
\]
se construye con las lecturas HMT ya fijadas de \(\pi\) y \(e\), y con los
enteros \(270\) y \(729\) del registro. Aquí es una entrada matemática
pre-dimensional: no intervienen unidades ni observables físicos.

La construcción interna del escalar no determina por sí sola qué operador
debe multiplicar. Para mantener visible esta premisa, declaramos:
\begin{description}
  \item[\HDef --- defecto de cambio de hoja.] La coordenada dodecafásica del defecto
  de cambio de hoja es \(\Delta_4P_3\).
\end{description}
\HAdd fijará la forma aditiva de la acción, pero no selecciona \HDef: para cualquier
\(c\ge0\), el mismo argumento admitiría el par
\((\Pi_\pi^{(468)},cP_3)\). Demostrar \HDef directamente desde una transición
total del núcleo topológico de fase enriquecido es una condición adicional.

Definimos la suma de Kronecker
\[
 \boxed{
 D_{\rm switch}^{\rm dod}
 =\Pi_\pi^{(468)}\otimes I_{11}
 +\Delta_4 I_{468}\otimes P_3.
 }
\]
Los dos sumandos son positivos, conmutan y actúan en factores diferentes.
Por tanto,
\[
 D_{\rm switch}^{\rm dod}\ge0.
\]
El entrelazador \(U_W\) del
\cref{thm:entrelazador-dodecafase-witt} produce la carta Witt:
\[
\begin{aligned}
 D_{\rm switch}^{\rm Witt}
 &=(I_{468}\otimes U_W)
 D_{\rm switch}^{\rm dod}
 (I_{468}\otimes U_W)^*\\
 &=\Pi_\pi^{(468)}\otimes I_{E_{30}}
 +\Delta_4I_{468}\otimes Q_3.
\end{aligned}
\]
Así, el mismo operador positivo se lee en la dodecafase y en el módulo de
incidencia.

La cadena isométrica orientada identifica además los sectores activos
\[
 I_{L_{\rm or}}\otimes P_G,\qquad P_3,\qquad Q_3.
\]
No se afirma una equivalencia de los complementos completos de los tres
módulos: \(W_{GK}\) es una equivalencia parcial de los rangos
seleccionados. Ésa es exactamente la información necesaria para identificar
la segunda contribución de la traza.

\begin{theorem}[Alcance exacto del entrelazamiento y de la inducción de acción]
\label{p3-20260826:thm:alcance-induccion-accion}
Sean \(W_{GK}\) y \(U_W\) las isometrías parciales del
\cref{thm:identificacion-isometrica-tres-proyectores}. Entonces:
\begin{enumerate}
  \item para todo \(c\geq0\),
  \[
   W_{GK}^{*}(cP_3)W_{GK}
   =c(I_{L_{\rm or}}\otimes P_G),
   \qquad
   U_W(cP_3)U_W^{*}=cQ_3;
  \]
  por tanto, el entrelazamiento transporta el coeficiente, pero no lo
  selecciona;
  \item si \(P_{\rm sgn}\) y \(P_{\rm std}\) son los proyectores centrales
  de la descomposición
  \[
   P_3V_{11}\big|_{H_*}
   \simeq \operatorname{sgn}_{H_*}\oplus\operatorname{std},
   \qquad
   \operatorname{rank}P_{\rm sgn}=1,
   \quad
   \operatorname{rank}P_{\rm std}=2,
  \]
  todo operador positivo autoadjunto soportado en \(P_3V_{11}\) y
  \(H_*\)-equivariante tiene la forma
  \[
   D_{a,b}=aP_{\rm sgn}+bP_{\rm std},
   \qquad a,b\geq0.
  \]
  Incluso imponer la traza de \(\Delta_4P_3\) deja la familia
  \[
   a+2b=3\Delta_4.
  \]
  La isotropía \(a=b\) no se deduce de la simetría residual \(S_3\);
  \item la probabilidad uniforme del catálogo,
  \(\mu_U(U)=1/468\), es la traza canónica del cociente finito y su
  elevación compensada es la determinada por \(J_F\). No es la
  medida imagen de la probabilidad uniforme sobre semillas, cuyos átomos son
  \(1/729,1/243,1/54\). En particular,
  \[
   \mu_U(\mathscr R_\pi^{\rm micro})=\frac5{468},
   \qquad
   F_\#\mu_{\rm seed}(\mathscr R_\pi^{\rm micro})=\frac7{729};
  \]
  \item la ley completa
  \[
   \delta_*=\frac5{468}+\frac3{11}\Delta_4,
   \qquad
   \lambda_*=1-\delta_*,
   \qquad
   a_*(e)=g(e)+\lambda_*s(e)
  \]
  es un teorema relativo a las cinco cláusulas \HTr--\HAdd--\HDef--\HRes--\HLoc. El motor
  ejecutable del núcleo topológico de fase reducido no induce por cubierta dirigida el operador ramificado
  \Gtr: la dinámica asociada retorna a la identidad tras \(54\) iteraciones,
  mientras cada estado de
  \Gtr posee cuatro sucesores distintos y ninguna autoarista.
\end{enumerate}
En consecuencia, los resultados anteriores demuestran la canonicidad de las identificaciones
isométricas, de la traza sobre el catálogo cociente y del carácter aditivo
dentro de sus premisas. La afirmación más fuerte de que una transición total
del núcleo topológico de fase enriquecido induce necesariamente \(\Delta_4P_3\), la medida uniforme y la
ley local requiere todavía el teorema dinámico adicional especificado en el
análisis de inducción.
\end{theorem}

\begin{proof}
Las identidades del primer punto se obtienen multiplicando
\(W_{GK}^{*}W_{GK}=I_{L_{\rm or}}\otimes P_G\),
\(W_{GK}W_{GK}^{*}=P_3\) y \(U_WP_3U_W^{*}=Q_3\) por el escalar \(c\).
Para el segundo punto, la restricción a \(H_*\simeq S_3\) es suma sin
multiplicidades de dos representaciones irreducibles no isomorfas. El lema de
Schur diagonaliza todo operador autoadjunto equivariante en esos dos bloques;
la positividad equivale a \(a,b\geq0\). Su traza normalizada es
\((a+2b)/11\), lo que prueba la libertad residual aun después de fijar la
traza. La acción completa de \(A_5\) sobre el triplete irreducible reduciría
la familia a \(cP_3\), pero tampoco fijaría el valor de \(c\).

El tercer punto sigue del censo de fibras: \(F\) tiene multiplicidades
\(144,432,1944\), de modo que la medida imagen pondera una emisión por su
número de presentaciones. La traza \(\tau_{468}\), en cambio, asigna el
mismo peso a cada clase; \(J_F\) realiza exactamente su elevación
normalizado. El cuarto punto es la concatenación lógica ya demostrada:
\HTr fija la traza de cociente, \HAdd la suma de trazas, \HDef declara
\(\Delta_4P_3\), \HRes define la capacidad residual y \HLoc la lleva a las aristas.
La enumeración finita prueba la obstrucción entre el núcleo reducido y
\Gtr: una cubierta dirigida preservaría la estrella saliente y una
semiconjugación del retorno identidad exigiría autoaristas en la imagen,
contradiciendo ambos censos de \Gtr.
\end{proof}

El espectro completo de \(D_{\rm switch}^{\rm dod}\) es
\[
\begin{array}{c|c}
\text{autovalor}&\text{multiplicidad}\\ \hline
0&(468-5)(11-3)=3704\\
\Delta_4&(468-5)3=1389\\
1&5(11-3)=40\\
1+\Delta_4&5\cdot3=15.
\end{array}
\]
Por consiguiente,
\[
\begin{aligned}
 \tau_{468\cdot11}(D_{\rm switch})
 &=\tau_{468}(\Pi_\pi^{(468)})+\Delta_4\tau_{11}(P_3)\\
 &=\boxed{\frac5{468}+\frac3{11}\Delta_4}.
\end{aligned}
\]
La misma identidad vale en la carta Witt. En la carta del sector orientado, el término
\(3/11\) es la traza de \(P_G\) sobre el espacio homogéneo de once ramas.

\section{Propiedad universal del funcional positivo de acción}

La palabra \emph{acción} se usa en sentido aditivo: la concatenación de dos
rutas recibe la suma de sus incrementos. Para separar este contenido de una
fórmula escogida, definimos primero la categoría en la que la ley será única.

\begin{definition}[Carácter positivo de acción HMT]
\label{p3-20260826:def:caracter-accion-hmt}
Un carácter local de acción sobre el producto
regional--dodecafásico es una aplicación
\[
 \delta:
 M_{468}(\mathbb C)_+\oplus M_{11}(\mathbb C)_+
 \longrightarrow\mathbb R_{\ge0}
\]
que satisface:
\begin{enumerate}
  \item aditividad exacta en el cono producto;
  \item homogeneidad positiva;
  \item invariancia por conjugación unitaria independiente en los dos
  factores;
  \item normalización bifactorial,
  \[
  \delta(I_{468},0)=\delta(0,I_{11})=1.
  \]
\end{enumerate}
\end{definition}

La aditividad traduce la concatenación; la invariancia elimina las
coordenadas; la normalización bifactorial fija una escala en cada sumando.
No se trata de un funcional unital sobre el álgebra suma directa: su unidad
es \((I_{468},I_{11})\) y la aditividad da
\(\delta(I_{468},I_{11})=2\). Ninguna cláusula contiene el valor de un
ángulo.

\begin{theorem}[Unicidad del carácter positivo y capacidad residual relativa]
\label{p3-20260826:thm:ley-accion-hmt}
En la clase de la definición anterior existe un único carácter positivo de
acción HMT\@. Es
\[
 \boxed{\delta(A,B)=\tau_{468}(A)+\tau_{11}(B).}
\]
Aplicado al defecto declarado por \HDef,
\((\Pi_\pi^{(468)},\Delta_4P_3)\), produce
\[
 \delta_*=\frac5{468}+\frac3{11}\Delta_4.
\]
Si, además, se declara \HRes, según la cual la capacidad autodual de cambio de
hoja se normaliza en uno y el defecto la reduce mediante
\(\lambda=1-\delta\), la capacidad residual dentro de \HTr--\HAdd--\HDef--\HRes es
\[
 \boxed{
 \lambda_*
 =1-\delta_*
 =1-\frac5{468}-\frac3{11}\Delta_4.
 }
\]
\end{theorem}

\begin{proof}
La aditividad separa los factores:
\[
 \delta(A,B)=\delta(A,0)+\delta(0,B).
\]
Cada restricción es positiva, homogénea e invariante por conjugación. Para
un proyector de rango uno, la invariancia unitaria da un valor común. La
aditividad sobre una descomposición espectral fuerza que cada restricción
sea un múltiplo de la traza. La normalización bifactorial fija esos múltiplos en
\(1/468\) y \(1/11\). De aquí resulta la suma de trazas. La última fórmula
aplica la normalización declarada de capacidad residual.
\end{proof}

Numéricamente,
\[
 \boxed{\lambda_*=0.9892859130191415\ldots},
 \qquad
 0<\lambda_*<1.
\]
La positividad se deduce ya de
\(0<\Delta_4<10^{-3}\) y \(5/468<1\). No depende de una comparación
metrológica.

\section{Functor de acción positiva sobre la categoría de caminos compatibles}

Recordemos el grafo dirigido \(\Gtr\) y su categoría libre
\(\mathsf P(\Gtr)\). Para una arista
\(e:x\to y\), definimos
\[
 g(e)=\mathbf1_{\{d_y\ne d_x\}},
 \qquad
 s(e)=\mathbf1_{\{m_y\ne m_x\}},
\]
donde \(g\) registra un giro y \(s\) un cambio de hoja. La flecha desde el
escalar \(\lambda_*\) hacia las aristas no se deduce del
teorema de trazas. Se declara mediante la cláusula \HLoc: un giro posee coste
unitario, un cambio de hoja posee coste \(\lambda_*\), y ambos indicadores se
suman cuando ocurren en la misma arista. Bajo \HLoc, la acción local es
\[
 \boxed{a_*(e)=g(e)+\lambda_*s(e).}
\]
Para una ruta \(\gamma=e_1\cdots e_n\),
\[
 \mathcal A_*(\gamma)
 =\sum_{j=1}^na_*(e_j)
 =\#\operatorname{giros}(\gamma)
 +\lambda_*\#\operatorname{cambios\ de\ hoja}(\gamma).
\]
Para la ruta vacía \(\mathbf1_x\), se fija
\(\mathcal A_*(\mathbf1_x)=0\). Entonces
\[
 \boxed{
 \mathcal A_*(\gamma_1\gamma_2)
 =\mathcal A_*(\gamma_1)+\mathcal A_*(\gamma_2).
 }
\]
Por tanto,
\[
 \mathcal A_*:\mathsf P(\Gtr)\longrightarrow(\mathbb R_{\ge0},+)
\]
es un funtor monoidal, o cociclo aditivo de caminos.

Los indicadores \(g\) y \(s\) son insensibles al sentido de recorrido. La
orientación no se pierde: reside en la línea \(L_{\rm or}\) y en la
involución de estrella, mientras el coste positivo reside en los
proyectores. Confundir ambos datos introduciría signos negativos en la
acción.

Para \(\beta>0\), el operador de transferencia asociado es
\[
 (\mathcal L_{\beta,*}f)(y)
 =\sum_{e:x\to y}e^{-\beta a_*(e)}f(x).
\]
Todos sus pesos son estrictamente positivos. Como el potencial depende sólo
de los dos indicadores locales, conmuta con la acción gruesa
\((\mathbb Z/3\mathbb Z)^3\) y respeta la descomposición en \(27\)
sectores de Bloch usada en la sección siguiente.

Una dinámica determinista finita no ponderada actúa por matrices de
permutación y tiene valores singulares iguales a uno. El cociclo
\(\mathcal A_*\) convierte el inventario de caminos en un operador de
transferencia que preserva el cono positivo, es decir, en una matriz con
pesos estrictamente positivos sobre las aristas permitidas. Esa propiedad no
demuestra por sí sola contracción, simplicidad del autovalor dominante ni
separación espectral; estas conclusiones requieren una normalización o un
argumento espectral adicional.

\ifdefined\HMTInternalTraceability
\section{Inducción por el cociente y extensión dinámica}

La frase \emph{inducida por el núcleo topológico de fase enriquecido} posee dos sentidos que deben mantenerse
separados.

\begin{enumerate}
  \item \textbf{Inducción por el objeto cociente.} El emisor ejecutable
  determina \(F\), su factor polar \(J_F\), el espacio \(H_U\) y la
  traza normalizada. El estado enriquecido aporta
  \(\Pi_\pi^{(468)}\), \(P_3\),
  \(K\), \(B_K\), el marco Witt y \(\Delta_4\). En este sentido
  categorial, \(D_{\rm switch}\) se construye sólo con datos internos
  APP--TRIT--TPK\@.
  \item \textbf{Inducción por una transición total.} Esta lectura exigiría
  una transición finita ejecutable sobre todas las coordenadas de
  \(\mathcal X_{\mathrm{TPK}}^{\mathrm{enr}}\) y una prueba de que su
  medida imagen dinámica
  preserva \(J_FH_U\) y produce el potencial \(a_*\). Esa transición total
  no pertenece al dominio ejecutable construido en esta sección.
\end{enumerate}

El \cref{p3-20260826:thm:ley-accion-hmt} y la construcción de caminos cierran la primera
lectura mediante cinco cláusulas explícitas del lector:
\begin{description}
  \item[\HTr --- traza de cociente.] Cuando la multiplicidad de presentación
  permanece en la fibra olvidada, una emisión distinta es un átomo y el
  estado del catálogo es la traza normalizada de su álgebra finita.
  \item[\HAdd --- cociclo de acción.] El defecto de una ruta es el carácter
  positivo, bifactorialmente normalizado, aditivo e invariante del producto
  regional--dodecafásico.
  \item[\HDef --- defecto de cambio de hoja.] El segundo argumento del carácter es
  \(\Delta_4P_3\).
  \item[\HRes --- capacidad residual.] La capacidad autodual se normaliza en
  uno y el defecto la reduce según \(\lambda=1-\delta\).
  \item[\HLoc --- acción local.] Sobre las aristas de \Gtr se fija
  \(a(e)=g(e)+\lambda s(e)\), y la acción de una ruta es la suma sobre sus
  aristas.
\end{description}
\HTr fuerza \(5/468\). \HAdd fuerza la suma de trazas y excluye términos
mixtos. \HDef fija el escalar que ocupa la coordenada de switch. \HRes transforma
el defecto en capacidad residual y \HLoc transporta esa capacidad a las aristas.
Así, \(\delta_*\) es consecuencia de \HTr--\HAdd--\HDef, \(\lambda_*\) lo es de
\HTr--\HAdd--\HDef--\HRes y \(\mathcal A_*\) lo es de
\HTr--\HAdd--\HDef--\HRes--\HLoc. \HTr y \HAdd poseen caracterizaciones universales; \HDef, \HRes y
\HLoc son identificaciones constitutivas internas; no se deducen de la
transición finita anterior.

\section{Familias compatibles con las simetrías parciales}

\subsection{Simetría de \(1/2\) ciclo y familia compatible de medidas}

El intercambio de las dos mitades de una emisión descompone las \(468\)
emisiones en \(234\) pares. Tanto la medida uniforme del catálogo como las
medidas imagen de semillas son invariantes bajo ese intercambio. También lo
es cualquier asignación constante en cada par. La invariancia del retorno,
por sí sola, no prueba uniformidad.

\subsection{Simetría de \texorpdfstring{\Gtr}{Gamma tr} y familia de costes}

Las \(7776\) aristas de \Gtr se dividen en \(288\) órbitas de tamaño
\(27\) bajo \((\mathbb Z/3\mathbb Z)^3\). Las \(2160\) aristas de
cambio de hoja forman \(80\) órbitas:
\[
 80=5\cdot16.
\]
La simetría permite asignar costes diferentes a esas órbitas. El coste
escalar común procede de \HAdd, no de la cardinalidad.

\subsection{Libertad escalar de la \(2^{\mathrm a}\) coordenada}

Para todo \(c\ge0\), el operador
\[
D_c=\Pi_\pi^{(468)}\otimes I_{11}
+cI_{468}\otimes P_3
\]
es positivo y el carácter único de \HAdd da
\[
\tau(D_c)=\frac5{468}+\frac3{11}c.
\]
La fórmula interna de \(\Delta_4\) construye un candidato distinguido, pero
la propiedad universal de la traza no demuestra por sí sola
\(c=\Delta_4\). Esa igualdad es precisamente el contenido de \HDef.

\subsection{Positividad, linealidad y términos de orden \(2\)}

Las funciones
\[
 1-u-v+cuv,\qquad(1-u)(1-v),\qquad e^{-(u+v)}
\]
pueden compartir normalización y primer orden. La ausencia de \(uv\) se
deduce de la aditividad exacta. Sin \HAdd, la ley lineal no sería única.

\section{Hipótesis de existencia, positividad, naturalidad y extensión dinámica}

El resultado depende de las siguientes condiciones verificables:
\begin{enumerate}
  \item el catálogo no contiene \(468\) emisiones o la microfibra no tiene
  rango cinco;
  \item el censo de fibras no es
  \(432\times144,\ 18\times432,\ 18\times1944\);
  \item \(J_F\) no es isométrico o no entrelaza observables del cociente;
  \item alguno de \(P_G,P_3,Q_3\) deja de ser un proyector de rango tres;
  \item existe otro carácter positivo, bifactorialmente normalizado, aditivo e
  invariante que no sea la suma de trazas;
  \item una transición total del núcleo topológico de fase enriquecido, una vez construida, induce un
  coeficiente de switch distinto de \(\Delta_4\);
  \item \(D_{\rm switch}\) adquiere un autovalor negativo;
  \item la acción acumulada deja de ser aditiva bajo concatenación.
\end{enumerate}

Los certificados reconstruyen las fibras, las dos medidas, el factor polar,
el espectro del defecto, las órbitas de aristas y las ablaciones. El valor de
\(\lambda_*\) aparece al final como consecuencia de esos objetos.
\fi
